\section{Theory}
\label{sec:theory}

Almost every attention-graph diagnostic begins by averaging the attention matrix over heads, and then reads a
spectral quantity off the single graph that remains \citep{noel2025gsp,li2025hsad,badash2026layers}. This section
shows that the averaging step, not the choice of spectral quantity, is what decides whether the diagnostic can
see a failure at all. Three statements follow. The first is about any functional of the averaged graph, the
second gives the direction and the equality condition of the error that averaging makes for algebraic
connectivity, and the third draws out what the features of the strongest published attention-only detector
measure, which is a degree sequence rather than a global spectrum.

\subsection{Setup}
\label{sec:theory:setup}

Fix one layer with $H$ heads and $T$ token positions. Head $h$ produces a row-stochastic, causal attention
matrix $A_h \in \mathbb{R}^{T \times T}$, with $(A_h)_{ij} = 0$ for $j > i$, which we symmetrise into a
weighted adjacency and strip of self-loops,
\begin{equation}
W_h = \tfrac{1}{2}\left(A_h + A_h^{\top}\right) - \operatorname{diag}\!\left(\tfrac{1}{2}(A_h + A_h^{\top})\right),
\label{eq:adj}
\end{equation}
with combinatorial Laplacian $L_h = D_h - W_h$, where $D_h = \operatorname{diag}(W_h \mathbf{1})$, and symmetric
normalised Laplacian $\mathcal{L}_h = I - D_h^{-1/2} W_h D_h^{-1/2}$ on the support of $D_h$. The
\emph{head-averaged} graph is $\bar W = \frac{1}{H}\sum_h W_h$, with $\bar L$ and $\bar{\mathcal{L}}$ built from
it. We write $\lamtwo(\cdot)$ for the second-smallest eigenvalue, the algebraic connectivity, and $\lammax$ for
the largest. A \emph{per-head} read of the layer keeps the vector
$\left(\phi(W_1),\dots,\phi(W_H)\right)$ for a scalar functional $\phi$; a \emph{head-averaged} read keeps the
single number $\phi(\bar W)$.

\subsection{Averaging is a lossy projection, whatever is read off it}
\label{sec:theory:annihilation}

\begin{proposition}[annihilation]
\label{prop:annihilation}
Let $\Phi$ be any function of the head-averaged graph, that is, any map of the form
$\{W_h\}_{h\le H} \mapsto \Phi(\bar W)$. Then $\Phi$ is constant on every set of head configurations sharing the
same mean. In particular, for every $H \ge 2$ and every $T \ge 3$ there exist configurations $\{W_h\}$ and
$\{W_h'\}$ with $\bar W = \bar W'$ whose per-head connectivities differ by
$\max_h \lamtwo(L_h) - \max_h \lamtwo(L_h') = \Omega(1)$, so no head-averaged read can separate them while a
per-head read separates them by a constant.
\end{proposition}

\begin{proof}
The first claim is immediate, since $\Phi$ factors through $\bar W$. For the second, note first which
symmetric graphs Eq.~\ref{eq:adj} can produce from a causal head. Since $(A_h)_{ij} = 0$ for $j > i$, the
off-diagonal entries are $W_{ij} = \tfrac12 (A_h)_{ij}$ for $i > j$, so a symmetric zero-diagonal $W$ is
realisable exactly when every node sends total weight at most $\tfrac12$ to the nodes before it,
$2\sum_{j<i} W_{ij} \le 1$ for all $i$: set $(A_h)_{ij} = 2W_{ij}$ for $i > j$ and put the remaining mass on
the diagonal. The condition is convex, so a mean of realisable graphs is realisable. Take $T \ge 3$, let $P$
be the path graph on $T$ nodes with weights $\tfrac12$, in which every node attends fully to its predecessor,
and let $K$ be the complete graph with weights $1/(2(T-1))$; both satisfy the condition. Put $W_1 = P$,
$W_2 = K$ and $W_1' = W_2' = \tfrac12(P+K)$, and for $H > 2$ repeat the pattern. Both configurations have
mean $\tfrac12(P+K)$. The path has $\lamtwo(L_P) = 1-\cos(\pi/T) = \Theta(T^{-2})$ while the complete graph
has $\lamtwo(L_K) = T/(2(T-1)) > \tfrac12$, so the per-head vectors differ in their maximum by more than
$\tfrac12 - \Theta(T^{-2})$, whereas every head-averaged functional returns the same value on both.
\end{proof}

Proposition~\ref{prop:annihilation} is elementary, and we state it because the literature it addresses has
behaved as though it were false: averaging is a many-to-one map, and the proposition only makes precise what
that map cannot preserve and shows the loss is of constant size on admissible graphs. It says nothing about
which spectral quantity is read, so it applies to the Fiedler value, to the whole eigenvalue
profile, to Dirichlet energy, and to any statistic of the averaged graph that a future paper might propose. The
information that distinguishes the two configurations is carried by the \emph{contrast between heads}, and the
mean is precisely the statistic that discards contrast. Whether real failures are visible in that contrast is an
empirical question, which Section~\ref{sec:exp:perhead} answers, but no averaged read can answer it either way.

\subsection{The direction of the error, and when it vanishes}
\label{sec:theory:jensen}

Proposition~\ref{prop:annihilation} exhibits a worst case. The next statement is about the typical case, and it
fixes the sign of the discrepancy.

\begin{proposition}[head averaging overstates connectivity]
\label{prop:jensen}
For any heads $W_1,\dots,W_H$,
\begin{equation}
\lamtwo(\bar L) \;\ge\; \frac{1}{H}\sum_{h=1}^{H} \lamtwo(L_h),
\label{eq:jensen}
\end{equation}
with equality if and only if the heads admit a common Fiedler vector, that is, a unit $x \perp \mathbf{1}$
attaining $\min_{y \perp \mathbf{1}, \|y\|=1} y^{\top} L_h y$ simultaneously for every $h$.
\end{proposition}

\begin{proof}
The Laplacian is linear in the adjacency, so $\bar L = \frac{1}{H}\sum_h L_h$. Since $L_h \mathbf{1} = 0$ and
$L_h \succeq 0$, Courant--Fischer gives
$\lamtwo(L) = \min_{x \perp \mathbf{1},\, \|x\|=1} x^{\top} L x$, a pointwise minimum of functions that are
linear in $L$. A pointwise minimum of linear functions is concave, so $\lamtwo$ is concave on the convex cone of
graph Laplacians, and Jensen's inequality applied to the uniform average over heads gives
Eq.~\ref{eq:jensen}. For the equality case, let $x^\star$ attain the minimum for $\bar L$. Then
$\lamtwo(\bar L) = \frac1H \sum_h (x^\star)^{\top} L_h x^\star \ge \frac1H \sum_h \lamtwo(L_h)$, and the
inequality is tight exactly when $(x^\star)^{\top} L_h x^\star = \lamtwo(L_h)$ for every $h$, that is, when
$x^\star$ is a Fiedler vector of every head. \end{proof}

The consequence is one-directional and therefore diagnostically bad in a specific way. A hallucinated call that
fragments the routing of a few heads lowers $\lamtwo$ for those heads, but Eq.~\ref{eq:jensen} bounds the
averaged graph's connectivity \emph{below} by the mean over heads, and the minimising direction of the average is
a consensus direction that the fragmented heads need not share. Averaging thus reports a graph that is at least
as well connected as the typical head, and the evidence of fragmentation is compressed toward the consensus.

The equality condition is a measure-zero event, so the question is not whether it holds but how far real layers
sit from it, and that is measurable. On
\nJensenLayers{} layers drawn from \nJensenModels{} models and three inputs, the combinatorial inequality
Eq.~\ref{eq:jensen} holds with \combGapViolations{} violations, as it must, and the gap is bounded away from zero
everywhere, with minimum \combGapMin{} and mean \combGapMean{} (Figure~\ref{fig:jensen}). The mean pairwise absolute
alignment between
per-head Fiedler vectors is \fiedlerAlignmentMean{}, far from the value $1$ that equality requires, and the
standard deviation of $\lamtwo$ across the heads of a layer averages \perHeadSdMean{}. For the normalised
Laplacian actually used by our detectors the Rayleigh quotient is not linear in $W$, so Eq.~\ref{eq:jensen} is
not a theorem; empirically it holds on all but \normGapViolations{} of the \nJensenLayers{} layers, and the
averaged graph reports connectivity a median \normRatioMedianPct{}\% above the mean over heads, reaching
$\normRatioMax{}\times$ at the extreme. Head averaging is therefore not a mild smoothing of a per-head quantity;
it reports a systematically better connected graph than any head in the layer, and it does so by an amount that
varies from layer to layer.

\begin{figure}[t]
\centering
\includegraphics[width=0.85\linewidth]{fig_jensen}
\caption{Proposition~\ref{prop:jensen} measured. Each point is one layer of one model on one input; the
horizontal axis is the mean over heads of the per-head algebraic connectivity and the vertical axis the
connectivity of the head-averaged graph. Every point lies above the diagonal for the combinatorial Laplacian,
as the proposition requires, and all but \normGapViolations{} do for the normalised Laplacian the detectors
use, where the inequality is not a theorem. The distance from the diagonal is the connectivity that head
averaging invents.}
\label{fig:jensen}
\end{figure}

\subsection{The strongest attention-only baseline reads degrees, not a spectrum}
\label{sec:theory:lapeig}

LapEigvals \citep{binkowski2025lapeigvals} is the strongest published attention-only hallucination detector. Its
features are the top-$k$ eigenvalues of the Laplacian $L = D - A$ of the unsymmetrised causal attention matrix,
and its authors observe (their Lemma~1) that because $A_{ij}=0$ for $j>i$ the matrix $L$ is lower triangular,
so its eigenvalues are its diagonal entries and no eigendecomposition is needed:
\begin{equation}
\operatorname{spec}(L) = \{\, d_j - A_{jj} \,\}_{j \le T},
\label{eq:lapeigdiag}
\end{equation}
where $d$ is their position-normalised in-degree. We take that observation as given and draw out what it
implies for the comparison in this paper.

\begin{remark}[LapEigvals is a degree statistic]
\label{prop:lapeig}
By Eq.~\ref{eq:lapeigdiag} the LapEigvals features depend on $A$ only through the column sums and the diagonal.
Any two causal attention matrices with the same in-degree sequence and the same self-attention weights have
identical features, however differently they route attention between positions. The features are therefore a
per-position, per-head degree statistic: local, invariant to in-degree-preserving rewiring, and carrying no
eigenvector and no connectivity information. The symmetrisation in Eq.~\ref{eq:adj} destroys triangularity,
so the spectra our detectors read are global functionals of the routing that change under such rewiring.
\end{remark}

This is not a criticism of the method, which works, and works well enough that our detectors do not beat it
(Section~\ref{sec:exp:frontier}). It is a statement about what the two families measure: they are not two
estimates of the same object, and the fact that both keep heads apart is the one thing they share.

\begin{corollary}[decorrelation prediction]
\label{cor:decorrelate}
If the two families measure different functionals, their scores should not be interchangeable: the rank
correlation between them should be well below one, each should retain predictive signal once the other is
conditioned on, and a combination of the two should exceed both. Section~\ref{sec:exp:decorrelate} tests the
three predictions; two hold and the third fails.
\end{corollary}

\subsection{Why the call-local subgraph is length-invariant}
\label{sec:theory:length}

The detectors read the subgraph induced on the generated call rather than on the whole sequence, which matters
for a confound. Lengths differ systematically between correct and hallucinated calls, so any feature that scales
with $T$ can carry label information that has nothing to do with routing. The normalised Laplacian removes the
crudest version of this.

\begin{lemma}[bounded spectrum]
\label{lem:bounded}
For any weighted graph with non-negative weights and no isolated vertices,
$\operatorname{spec}(\mathcal{L}) \subset [0,2]$, with $\lammax = 2$ if and only if a connected component is
bipartite. The bound is independent of the number of nodes.
\end{lemma}

\begin{proof}
Standard \citep{chung1997spectral}: for $x \ne 0$,
$\frac{x^\top \mathcal{L} x}{x^\top x}$ equals a ratio whose numerator
$\sum_{i<j} W_{ij} (y_i - y_j)^2$ with $y = D^{-1/2}x$ is bounded by
$2\sum_{i<j} W_{ij}(y_i^2 + y_j^2) = 2 x^\top x$, with equality only in the bipartite case. \end{proof}

Every per-head feature we use is a function of a spectrum contained in $[0,2]$ whatever the length of the call,
whereas the combinatorial spectrum grows with degree and hence with $T$. Length invariance by construction does
not prove the absence of a length confound, since the routing itself may correlate with length, so
Section~\ref{sec:exp:confound} measures a surface-feature baseline built from lengths alone and reports every
detector against it.
